\section{Final Verdict on Industrial Feasibility}

The deployment of the proposed dual-arm collaborative cell represents a highly feasible and advantageous transition for medium-scale, high-mix assembly operations. The analysis of this architecture across economic, technical, and safety domains confirms its readiness for industrial application, provided that specific integration strategies are employed.

From an economic perspective, the initial Capital Expenditure (CAPEX) required for the dual-manipulator hardware, advanced vision systems, and edge-computing infrastructure is heavily offset by the rapid Return on Investment (ROI). Industry data establishes that collaborative setups typically achieve payback within 12 to 18 months \cite{ifactory2026calculator, mantec2026cobot}. This is driven by a 20\% to 40\% increase in throughput, enabled by the multithreaded parallel execution framework \cite{szghtech2026cobot}, and a significant reduction in material scrap due to the predictive accuracy of the Digital Twin \cite{iotforall2026digitaltwin}.

Technically, the transition from a laboratory prototype to a continuous production environment necessitates a hybrid control architecture. While the ROS 2 and MoveIt 2 stack provides unparalleled dynamic intelligence for complex task planning and Digital Twin synchronization, it must be paired with industrial PLCs to guarantee microsecond determinism and hardware-level safety interrupts \cite{mdpi2026interoperability}. This bifurcation ensures that the system benefits from advanced algorithmic flexibility without compromising industrial reliability.

Crucially, the system design aligns with stringent international safety regulations. The integration of the \texttt{ZoneDetectionManager} serves as a foundational software mechanism to satisfy the Speed and Separation Monitoring (SSM) requirements mandated by ISO/TS 15066 \cite{iso15066, nist2016ssm}. By proactively evaluating Cartesian proximity, the system eliminates the severe mechanical risks traditionally associated with multi-arm interference.

Ultimately, this architecture successfully resolves the automation bottleneck inherent in rigid, single-arm cells. By bridging the high-speed capability of an industrial backbone with the adaptability of a secondary collaborative arm, and safeguarding the physical execution through a real-time Digital Twin, this RRC system offers a commercially viable, scalable solution for the next generation of adaptive manufacturing.